\chapter*{前言}
\addcontentsline{toc}{chapter}{前言}

第三册原理卷把 AI 推理系统拆成模型资产、typed graph、张量布局、量化、KV cache、后端、服务运行时和验收体系。本实践与代码卷只做一件事：把这些概念压到同一个可运行项目里，并把完整源码阅读路径并入同一目录，让读者知道每个概念对应哪段代码、哪条命令、哪种输入、哪份报告和哪条失败路径。

贯穿项目是 \filepath{books/cpu-volume-3/labs/linux\_cpu\_inference}。它不是玩具目录，而是第三册已经持续维护的本地 CPU 量化推理切片：包含教学模型格式、tokenizer、safetensors manifest gate、int8 linear、packed/AVX2 路径、tiny reference decoder、GPT-2 reference path、KV cache trace、7B compute ledger、serving SLO probe、自研 GPT-2 smoke、真实 small 开源模型 smoke、CLI、benchmark 和 CTest。本卷不维护第二份实现；源码章节通过 \code{lstinputlisting} 直接读取同一份工程文件，所有练习、讲解和完整代码都回到同一份源码。

本卷的学习顺序是从小证据到大系统：先让一个 tiny MLP 算对，再解释张量地址流和 int8 累加；再让 reference decoder、sampler 和 KV cache trace 固定语义；然后接模型导入、shape verifier、memory planner 和 IR lowering；最后用 CPU/GPU 后端边界、工业专项和回归门禁收束。代码走读不再被放到最后当附录，而是插入到相同主题部分中：讲完项目入口就读构建文件和调用链，讲完量化就读公共合同、tiny MLP 和 int8 kernel，讲到真实模型时读 tokenizer、safetensors 和 GPT-2，收束验收时读工具、benchmark、smoke 和测试。每章都要求读者交付代码运行结果或报告字段，而不是只抄概念定义。

\begin{keyidea}
第三册实践的目标不是“写一个看起来像大模型的 demo”，而是建立一条能被复查的推理证据链：资产能校验，shape 能验证，kernel 能对照，trace 能定位，性能能降级，服务指标能解释，回归门禁能长期维护。
\end{keyidea}
